\documentclass[10pt,a4paper]{amsart}
\usepackage[margin=19mm]{geometry}
\usepackage{amsmath,amssymb,mathtools}
\usepackage[scr=rsfs,cal=euler]{mathalfa}
\usepackage{xfrac}
\usepackage[unicode,hidelinks,bookmarks=false]{hyperref}
\newcommand{\ie}{i.e.\ }
\newcommand{\cf}{cf.\ }
\newcommand{\resp}{respectively\ }
\newcommand{\Subsection}{\subsection}
\providecommand{\nobreakdash}{\nobreak\hbox{-}}
\setlength{\emergencystretch}{2em}
\multlinegap0pt
\usepackage[shortlabels]{enumitem}
\usepackage{mathtools}
\usepackage{xcolor}
\usepackage{float}
\usepackage{amssymb}
\usepackage{bm}
\usepackage[normalem]{ulem}
\usepackage{tensor}
\usepackage{tikz}
\usepackage{csquotes}
\newtheorem{proposition}{Proposition}
\newcommand{\Irr}{\text{Irr}}
\newcommand{\be}{\begin{equation}}
\newcommand{\cC}{\mathcal{C}}
\newcommand{\cL}{\mathcal{L}}
\newcommand{\cP}{\mathcal{P}}
\newcommand{\cS}{\mathcal{S}}
\newcommand{\cZ}{\mathcal{Z}}
\newcommand{\ee}{\end{equation}}
\title{Twin Algebras: \\ Condensable Algebras beyond Anyons}
\author{Yuhan Gai, Sakura Sch\"afer-Nameki, Alison Warman}
\date{Source: arXiv:2605.31602v1 (CC BY 4.0)}
\begin{document}
\maketitle

\noindent\emph{Source and licence.} Twin Algebras: \\ Condensable Algebras beyond Anyons. Version arXiv:2605.31602v1 (CC BY 4.0), distributed by arXiv under the Creative Commons Attribution 4.0 International licence, http://creativecommons.org/licenses/by/4.0/, as recorded in the arXiv licence metadata for this exact version and shown on the abstract page https://arxiv.org/abs/2605.31602v1. Full licence notice: \texttt{licenses/arxiv-2605.31602-CC-BY-4.0.txt}.\par\smallskip

\noindent\textit{Editorial scope.} Counted unit: the article's proposition prop:no-hidden-SB (source0.tex lines 1510--1516). The article's complete original proof of it is delivered with it (source0.tex lines 1517--1528, one \texttt{proof} environment of 12 lines). No mathematics is written, repaired or re-derived: the statement and the proof are the article's own lines, in the article's own notation, and no retained line is substituted or re-flowed.  No further source-local context is needed: every cross-reference of every retained line resolves inside the two delivered spans.  Nothing was omitted from any retained span, so the package carries no omission record.  No retained line contains a citation, so the wrapper carries no bibliography and no bibliography entry.  No retained line contains a figure environment, an image-inclusion command or drawing code, so no image is delivered and the package declares no asset dependency.  Editorial formatting only, in addition: an AMS \texttt{amsart} preamble that restates the article's own declarations and macros used by the retained lines, namely the environments \texttt{proposition} on separate counters, together with the standard packages those lines need.  The retained environments print in this document as Proposition 1 in order of appearance, whereas the article prints the same items with the numbering of its own full document.  The complete native source of the article as distributed in the arXiv e-print for this exact version is bundled unchanged as \texttt{sources/201736/source0.tex}.
\begin{proposition}[No relative SSBing]
\label{prop:no-hidden-SB}
    Let $\cL_1$ and $\cL_2$ be Lagrangian algebras in $\cZ(\cC)$. $\cP^{\cC}_{\cL_1}$ and $\cP^{\cC}_{\cL_2}$ have no relative SSBing
    if and only if $\cL_1 \cong \cL_2$ as objects in $\cZ(\cC)$, i.e. they are twin Lagrangian algebras.
    
    In particular, for any symmetry $\cS$ that is Morita equivalent to $\cC$ the phases $\cP^{\cS}_{\cL_1}$ and $\cP^{\cS}_{\cL_2}$ have the same symmetry breaking pattern. 
\end{proposition}

\begin{proof}
    Assuming $\cL_1\not\cong \cL_2$ as objects in $\cZ(\cC)$, then there exists $a\in \cL_1$ such that 
    \be
    n_{\cL_1,a}\neq n_{\cL_2,a}.
    \ee
    Let $\cC_1$ denote the fusion category corresponding to $\cL_1$, then the two subspaces 
    \be
        (V_{\cL_1}^{\cC_1})_{a,1} \not\cong (V_{\cL_2}^{\cC_1})_{a,1}
    \ee
    are inequivalent, showing that $\cP^{\cC_1}_{\cL_1}$ and $\cP^{\cC_1}_{\cL_2}$ admit distinct symmetry breaking patterns. We conclude either $\cP^{\cC}_{\cL_1}$ and $\cP^{\cC}_{\cL_2}$ have distinct symmetry breaking pattern, or if they have the same symmetry breaking pattern, they can be gauged to $\cP^{\cC_1}_{\cL_1}$ and $\cP^{\cC_1}_{\cL_2}$ with different symmetry breaking pattern, hence admitting relative SSBing. 
    For any $\cS$ symmetry dual to $\cC$ and $a\in \Irr(\cZ(\cC))$, $\cP^{\cS}_{\cL_1}$ and $\cP^{\cS}_{\cL_2}$ have isomorphic spaces $(V_{\cL_1}^{{\cS}})_{a,1}\cong (V_{\cL_2}^{{\cS}})_{a,1}$ of local operators, with respect to the $\cS$ symmetry. Therefore there is no relative SSBing between $\cP^{\cC_1}_{\cL_1}$ and $\cP^{\cC_1}_{\cL_2}$ as they admit the same symmetry breaking pattern for all symmetry categories related to $\cC$ by gauging. 
\end{proof}

\end{document}
